\section{შესავალი}

პროექტის თემის შერჩევა საკმაოდ პრაქტიკული დაკვირვებით დაიწყო. სწავლის
პერიოდში არაერთხელ აღმოვჩენილვართ სიტუაციაში, როდესაც სემინარზე ან
ლაბორატორიულ სამუშაოზე საჭირო იყო, რომ რამდენიმე სტუდენტს, ან სტუდენტსა და ლექტორს, ერთდროულად
ემუშავა ერთსა და იმავე კოდზე. სტანდარტული გამოსავალი, ეკრანის გაზიარება -
უნივერსიტეტის ქსელში ხშირად უბრალოდ არ მუშაობდა: ინტერნეტი წყდებოდა,
ვიდეო ნაკადი ჭედავდა და საბოლოოდ მთელი ჯგუფი ლექტორის ლეპტოპის გარშემო
იკრიბებოდა, რომ სემინარში მონაწილეობა მიეღო. Google Docs-ის ტიპის
ინსტრუმენტები კი, გარდა იმისა, რომ კოდის რედაქტირებისთვის მოუხერხებელია,
სტაბილურ ინტერნეტსა და ცენტრალურ სერვერს საჭიროებს, ანუ ზუსტად იმას,
რაც აუდიტორიაში ხშირად არ გვქონდა.

აქედან გაჩნდა კითხვა: შესაძლებელია თუ არა ისეთი ერთობლივი რედაქტორის
შექმნა, რომელიც ცენტრალურ სერვისზე საერთოდ არ იქნება დამოკიდებული და
იმუშავებს როგორც ლოკალურ ქსელში (თუნდაც ინტერნეტის გარეშე), ისე
ინტერნეტით დაშორებულ მონაწილეებთანაც და ექნება ownership სისტემა, 
რომელიც ლექტორს მისცემს საშუალებას, აკონტროლოს სესია და მართოს სტუდენტთა პრივილეგიები.
თემის ამ მიმართულებით განვითარებაში
დაგვეხმარა ჩვენი ხელმძღვანელი, ლუკა აბსანძე, რომელმაც პროექტის საწყის
ეტაპზე რამდენიმე საინტერესო იდეა და მიმართულება შემოგვთავაზა, საიდანაც
საბოლოო არჩევანი გავაკეთეთ.

პრობლემა, რომელსაც პროექტი ეხმიანება, შემდეგნაირად შეიძლება ჩამოყალიბდეს:
რეალურ დროში ერთობლივი ტექსტური რედაქტირების არსებული გადაწყვეტილებები
თითქმის ყველა შემთხვევაში მოითხოვს (1) დოკუმენტის მესამე მხარის სერვერზე
შენახვას, (2) მუდმივ ინტერნეტკავშირს და (3) ხშირად ანგარიშების შექმნასა და
კონკრეტულ ეკოსისტემაზე მიბმას. ეს სამივე მოთხოვნა პრობლემურია მაშინ,
როდესაც სამუშაო ფაილი კონფიდენციალურია, ქსელი არასტაბილურია ან უბრალოდ
გვინდა, რომ ერთი ოთახის ფარგლებში, ზედმეტი ინფრასტრუქტურის გარეშე
ვითანამშრომლოთ.

პროექტის მიზანი იყო სრულფასოვანი, კროსპლატფორმული desktop აპლიკაციის
შექმნა, რომელიც:
\begin{itemize}[itemsep=0pt]
  \item საშუალებას აძლევს ერთ მომხმარებელს (ჰოსტს), წამოიწყოს სესია და
        დანარჩენებს, უბრალოდ ბმულით შეუერთდნენ მას;
  \item პარალელურ რედაქტირებას უკონფლიქტოდ წყვეტს CRDT ალგორითმით, ისე,
        რომ ყველა მონაწილე გარანტირებულად ერთსა და იმავე ტექსტს ხედავდეს;
  \item მუშაობს როგორც ლოკალურ ქსელში ინტერნეტის გარეშე, ისე გლობალურად -
        Cloudflare Tunnel-ის მეშვეობით, ქსელის ყოველგვარი კონფიგურაციის
        გარეშე;
  \item ჰოსტს აძლევს სესიის სრულ კონტროლს: სტუმრების ჩაწერის უფლებების
        მართვას, სესიის დაწყება/დასრულებას და ფაილების ლოკალურად შენახვას.
\end{itemize}

პროექტის ფარგლები შეგნებულად შემოვსაზღვრეთ: სესია ერთ დოკუმენტზე მუშაობს,
დოკუმენტი უბრალო ტექსტია (კოდი, კონფიგურაცია, ჩანაწერები) და არა
ფორმატირებული rich text, ხოლო სესია ეფემერულია, ის ჰოსტის აპლიკაციასთან
ერთად იბადება და კვდება, თუმცა ნამუშევრის ფაილად შენახვა ნებისმიერ
მომენტშია შესაძლებელი.

მეთოდოლოგიის მხრივ სისტემა სამ დამოუკიდებლად აწყობად კომპონენტად დავყავით:
CRDT ბიბლიოთეკა Rust-ზე (\texttt{crdt-core}), desktop კლიენტი Tauri
ფრეიმვორკზე (Rust ბექენდი + React/TypeScript ფრონტენდი Monaco რედაქტორით)
და მსუბუქი relay სერვერი Go-ზე (\texttt{gateway}). 

პოტენციური მომხმარებლები, პირველ რიგში, სტუდენტები და ლექტორები არიან -
სემინარები, ლაბორატორიული სამუშაოები, ჯგუფური პროექტები, თუმცა იგივე
სცენარები ბუნებრივად ვრცელდება წყვილში პროგრამირებაზე, ტექნიკურ
გასაუბრებებზე, მენტორინგსა და ნებისმიერ სიტუაციაზე, სადაც რამდენიმე
ადამიანს ერთი ტექსტის ერთდროული რედაქტირება სჭირდება ცენტრალური
სერვისისთვის მონაცემების გადაცემის გარეშე.

და ბოლოს, პროექტს აქვს რამდენიმე მოსალოდნელი შეზღუდვა. სესია ჰოსტის
კომპიუტერზეა მიბმული: თუ ჰოსტი გაითიშა, სესია სრულდება (თუმცა თითოეულ
მონაწილეს დოკუმენტის სრული ასლი რჩება). გლობალური წვდომა დამოკიდებულია
Cloudflare-ის უფასო \texttt{trycloudflare.com} სერვისის ხელმისაწვდომობაზე -
მისი გაუმართაობის შემთხვევაში აპლიკაცია ლოკალური ქსელის რეჟიმში აგრძელებს
მუშაობას.
